% The bench as it stands today, and the parts that must arrive before chapter 9.
% Nothing is wired in this chapter. Positions are explicit so that the cable
% labels keep their clearance if the text is edited.
\begin{tikzpicture}[font=\sffamily\small, x=1cm, y=1cm]

% ---------------- the board
\node[board, minimum width=48mm, minimum height=30mm] (nucleo) at (0,1.2) {};
\node[text=white, font=\sffamily\bfseries\small] at (0,1.55) {NUCLEO-H7A3ZI-Q};
\node[text=white, font=\sffamily\scriptsize] at (0,1.15) {STM32H7A3ZI};

% the probe half of the board, and the empty shield footprint
\draw[draw=white, dashed, line width=0.5pt] (-0.85,-0.3) -- (-0.85,2.7);
\node[text=white, font=\sffamily\tiny, anchor=north, rotate=90] at (-1.05,1.9) {probe};
\node[pinrow, minimum width=42mm] at (0,2.45) {};
\node[chlabel, anchor=south] at (0,2.85) {Arduino header, empty until chapter 6};

% indicators, inside the board with their pins
\node[ledsym, fill=green!65!black]  (ld1) at (-0.4,0.15) {};
\node[ledsym, fill=yellow!80!black] (ld2) at ( 0.2,0.15) {};
\node[ledsym, fill=red!75!black]    (ld3) at ( 0.8,0.15) {};
\node[text=white, font=\sffamily\tiny] at (-0.4,0.55) {PB0};
\node[text=white, font=\sffamily\tiny] at ( 0.2,0.55) {PE1};
\node[text=white, font=\sffamily\tiny] at ( 0.8,0.55) {PB14};
\node[text=white, font=\sffamily\tiny] at (0.2,-0.13) {node state, chapter 1};

% ---------------- the host side
\node[ext, text width=26mm] (host) at (-6.9,1.2) {host PC\\{\scriptsize build, flash, terminal}};
\node[ext, text width=26mm] (pi)   at (-6.9,-0.9) {Raspberry Pi 4\\{\scriptsize the clean-clone build, today}};
\node[chlabel, below=1mm of pi] {the motion master from chapter 10};
\draw[usbcable] (host.east) -- node[lbl, above, align=center]{USB, one cable\\{\tiny power, SWD, console}} (nucleo.west);
\draw[bus, draw=linecol!60] (host) -- node[lbl, right]{network} (pi);

% ---------------- what has to be bought, one column
\node[absentblk, text width=32mm] (xcvr) at (6.3,2.6)
  {CAN-FD transceiver module\\{\scriptsize read the suffix: one letter separates
   the 1 Mbit part from the 5 Mbit part}};
\node[absentblk, text width=32mm] (term) at (6.3,0.3)
  {two 120 ohm terminators\\{\scriptsize one at each end of the bus, not one per node}};
\node[absentblk, text width=32mm] (adapter) at (6.3,-1.9)
  {host bus adapter\\{\scriptsize it must present the kernel's own socket
   interface, not a vendor library}};
\begin{scope}[on background layer]
\node[layer, fit=(xcvr)(term)(adapter), inner sep=6pt,
      label={[layerlabel, anchor=south west]north west:not yet on the bench}] {};
\end{scope}

\draw[hiz, ->] (nucleo.east) -- node[lbl, above]{PD0, PD1}
                                node[lbl, below]{chapter 9} (xcvr.west);

\node[umlnote, text width=62mm, anchor=north] at (-1.0,-2.2)
  {Nothing is wired in this chapter. The board carries no bus transceiver, which
   is the whole electrical gap between this node and a network. A 1 Mbit
   transceiver still gives full 64-byte frames and the stronger checksum,
   because the frame format is invisible to it; what it costs is only the faster
   data phase. The pin pair is confirmed in chapter 9 against the board manual,
   and the adapter is the purchase to get right: one inexpensive analyser ships
   a closed library with a frame structure of its own, which the standard
   command-line tools cannot see at all.};
\end{tikzpicture}
